% components/theorems.tex
% Charlie's signature mathematical environments, boxed theorems, ribbons, and shortcuts.

\RequirePackage[most,many,breakable]{tcolorbox}
\RequirePackage{varwidth}
\RequirePackage{tikz}
\usetikzlibrary{arrows,calc,shadows.blur}
\tcbuselibrary{theorems,skins,hooks}

\makeatletter

%==================================================
% TColorBox Styles
%==================================================

\tcbset{
    cmtleftpanel/.style 2 args={
        enhanced,
        breakable,
        colback = #1,
        frame hidden,
        boxrule = 0sp,
        borderline west = {2.5pt}{0pt}{#2},
        sharp corners,
        description font = \mdseries,
        separator sign none,
    },
    cmtleftthm/.style 2 args={
        cmtleftpanel={#1}{#2},
        detach title,
        before upper = \tcbtitle\par\smallskip,
        coltitle = #2,
        fonttitle = \bfseries\sffamily,
        segmentation style={solid, #2},
    },
    cmtexamplethm/.style={
        colback = myexamplebg,
        breakable,
        colframe = myexamplefr,
        coltitle = myexampleti,
        boxrule = 1pt,
        sharp corners,
        detach title,
        before upper=\tcbtitle\par\smallskip,
        fonttitle = \bfseries\sffamily,
        description font = \mdseries,
        separator sign none,
        description delimiters parenthesis,
    },
    cmtdefinitionthm/.style={
        enhanced,
        before skip=3mm,
        after skip=3mm,
        colback=mydefinitbg,
        colframe=mydefinitfr,
        boxrule=0.5mm,
        attach boxed title to top left={xshift=1cm,yshift*=1mm-\tcboxedtitleheight},
        varwidth boxed title*=-3cm,
        colbacktitle=mydefinitfr,
        coltitle=white,
        boxed title style={
            frame code={
                \path[fill=tcbcolback]
                ([yshift=-1mm,xshift=-1mm]frame.north west)
                arc[start angle=0,end angle=180,radius=1mm]
                ([yshift=-1mm,xshift=1mm]frame.north east)
                arc[start angle=180,end angle=0,radius=1mm];
                \path[left color=tcbcolback!60!black,right color=tcbcolback!60!black,
                    middle color=tcbcolback!80!black]
                ([xshift=-2mm]frame.north west) -- ([xshift=2mm]frame.north east)
                [rounded corners=1mm]-- ([xshift=1mm,yshift=-1mm]frame.north east)
                -- (frame.south east) -- (frame.south west)
                -- ([xshift=-1mm,yshift=-1mm]frame.north west)
                [sharp corners]-- cycle;
            },
            interior engine=empty,
        },
        fonttitle=\bfseries\sffamily,
    },
    cmtribbonbox/.style={
        enhanced,
        breakable,
        colback=white,
        attach boxed title to top left={yshift*=-\tcboxedtitleheight},
        fonttitle=\bfseries\sffamily,
        boxed title size=title,
        boxed title style={
            sharp corners,
            rounded corners=northwest,
            colback=tcbcolframe,
            boxrule=0pt,
        },
        underlay boxed title={
            \path[fill=tcbcolframe] (title.south west)--(title.south east)
            to[out=0, in=180] ([xshift=5mm]title.east)--
            (title.center-|frame.east)
            [rounded corners=\kvtcb@arc] |-
            (frame.north) -| cycle;
        },
    },
}

%==================================================
% Theorem Environments
%==================================================

% 1. Theorem
\newtcbtheorem[
    number within=chapter,
    crefname={theorem}{theorems},
    Crefname={Theorem}{Theorems}
]{Theorem}{Theorem}{cmtleftthm={mytheorembg}{mytheoremfr}}{th}

\newtcbtheorem[
    number within=chapter,
    crefname={theorem}{theorems},
    Crefname={Theorem}{Theorems}
]{theorem}{Theorem}{cmtleftthm={mytheorembg}{mytheoremfr}}{th}

\newtcolorbox{Theoremcon}{
    cmtleftpanel={mytheorembg}{mytheoremfr}
}

% 2. Corollary
\newtcbtheorem[
    number within=chapter,
    crefname={corollary}{corollaries},
    Crefname={Corollary}{Corollaries}
]{Corollary}{Corollary}{cmtleftthm={myp!10}{myp!85!black}}{th}

\newtcbtheorem[
    number within=chapter,
    crefname={corollary}{corollaries},
    Crefname={Corollary}{Corollaries}
]{corollary}{Corollary}{cmtleftthm={myp!10}{myp!85!black}}{th}

% 3. Lemma
\newtcbtheorem[
    number within=chapter,
    crefname={lemma}{lemmas},
    Crefname={Lemma}{Lemmas}
]{Lemma}{Lemma}{cmtleftthm={mylemmabg}{mylemmafr}}{th}

\newtcbtheorem[
    number within=chapter,
    crefname={lemma}{lemmas},
    Crefname={Lemma}{Lemmas}
]{lemma}{Lemma}{cmtleftthm={mylemmabg}{mylemmafr}}{th}

% 4. Proposition
\newtcbtheorem[
    number within=chapter,
    crefname={proposition}{propositions},
    Crefname={Proposition}{Propositions}
]{Prop}{Proposition}{cmtleftthm={mypropbg}{mypropfr}}{th}

\newtcbtheorem[
    number within=chapter,
    crefname={proposition}{propositions},
    Crefname={Proposition}{Propositions}
]{prop}{Proposition}{cmtleftthm={mypropbg}{mypropfr}}{th}

% 5. Claim
\newtcbtheorem[
    number within=chapter,
    crefname={claim}{claims},
    Crefname={Claim}{Claims}
]{claim}{Claim}{
    cmtleftthm={myg!10}{myg!85!black},
    borderline west = {2.5pt}{0pt}{myg},
}{th}

% 6. Exercise
\newtcbtheorem[
    number within=chapter,
    crefname={exercise}{exercises},
    Crefname={Exercise}{Exercises}
]{Exercise}{Exercise}{cmtleftthm={myexercisebg}{myexercisefg}}{th}

\newtcbtheorem[
    number within=chapter,
    crefname={exercise}{exercises},
    Crefname={Exercise}{Exercises}
]{exercise}{Exercise}{cmtleftthm={myexercisebg}{myexercisefg}}{th}

% 7. Example
\newtcbtheorem[
    number within=chapter,
    crefname={example}{examples},
    Crefname={Example}{Examples}
]{Example}{Example}{cmtexamplethm}{ex}

\newtcbtheorem[
    number within=chapter,
    crefname={example}{examples},
    Crefname={Example}{Examples}
]{example}{Example}{cmtexamplethm}{ex}

% 8. Definition (Signature Bookmark Tab Style)
\newtcbtheorem[
    number within=chapter,
    crefname={definition}{definitions},
    Crefname={Definition}{Definitions}
]{Definition}{Definition}{
    cmtdefinitionthm,
    title={#2},#1
}{def}

\newtcbtheorem[
    number within=chapter,
    crefname={definition}{definitions},
    Crefname={Definition}{Definitions}
]{definition}{Definition}{
    cmtdefinitionthm,
    title={#2},#1
}{def}

% 9. Remark
\newtcbtheorem[
    number within=chapter,
    crefname={remark}{remarks},
    Crefname={Remark}{Remarks}
]{Remark}{Remark}{cmtleftthm={myremarkbg}{myremarkfr}}{rmk}

\newtcbtheorem[
    number within=chapter,
    crefname={remark}{remarks},
    Crefname={Remark}{Remarks}
]{remark}{Remark}{cmtleftthm={myremarkbg}{myremarkfr}}{rmk}

% 10. Question (Signature Ribbon Style)
\@ifundefined{c@questioncounter}{\newcounter{questioncounter}}{}
\@ifundefined{c@chapter}{}{%
    \counterwithin{questioncounter}{chapter}%
}

\newtcbtheorem[
    use counter=questioncounter,
    crefname={question}{questions},
    Crefname={Question}{Questions}
]{question}{Question}{
    cmtribbonbox,
    colframe=myb!80!black,
    title={#2},
    #1
}{qs}

\newtcbtheorem[
    crefname={question}{questions},
    Crefname={Question}{Questions}
]{qstion}{Question}{
    cmtribbonbox,
    colframe=mygr,
    title={#2},
    #1
}{qs}

% 11. Solution (Ribbon Style)
\newtcolorbox{solution}[1][]{
    cmtribbonbox,
    colframe=myg!80!black,
    title=Solution,
    #1
}
\newenvironment{solbox}[1][]{\begin{solution}[#1]}{\end{solution}}

% 12. Wrong Concept Box (Red Top Badge)
\newtcbtheorem[
    number within=chapter,
    crefname={wrong concept}{wrong concepts},
    Crefname={Wrong Concept}{Wrong Concepts}
]{wconc}{Wrong Concept}{
    breakable,
    enhanced,
    colback=white,
    colframe=myr,
    arc=0pt,
    outer arc=0pt,
    fonttitle=\bfseries\sffamily\large,
    colbacktitle=myr,
    attach boxed title to top left={},
    boxed title style={
        enhanced,
        skin=enhancedfirst jigsaw,
        arc=3pt,
        bottom=0pt,
        interior style={fill=myr}
    },
    #1
}{wc}

% 13. Note Box (Card with Rivet / Screw Holes & Blur Drop Shadow)
\newtcolorbox{note}[1][]{
    enhanced jigsaw,
    colback=gray!15!white,
    colframe=gray!80!black,
    size=small,
    boxrule=1pt,
    title=\textbf{\sffamily Note:},
    halign title=flush center,
    coltitle=black,
    breakable,
    drop shadow=black!40!white,
    attach boxed title to top left={xshift=1cm,yshift=-\tcboxedtitleheight/2,yshifttext=-\tcboxedtitleheight/2},
    minipage boxed title=1.8cm,
    boxed title style={
        colback=white,
        size=fbox,
        boxrule=1pt,
        boxsep=2pt,
        underlay={
            \coordinate (dotA) at ($(interior.west) + (-0.5pt,0)$);
            \coordinate (dotB) at ($(interior.east) + (0.5pt,0)$);
            \begin{scope}
                \clip (interior.north west) rectangle ([xshift=3ex]interior.east);
                \filldraw [white, blur shadow={shadow opacity=60, shadow yshift=-.75ex}, rounded corners=2pt] (interior.north west) rectangle (interior.south east);
            \end{scope}
            \begin{scope}[gray!80!black]
                \fill (dotA) circle (2pt);
                \fill (dotB) circle (2pt);
            \end{scope}
        },
    },
    #1,
}

% 14. Proof Environment
\newenvironment{myproof}[1][\proofname]{%
    \begin{proof}[\bfseries\sffamily #1]%
}{%
    \end{proof}%
}

\newenvironment{myclaim}[1][Claim]{\begin{proof}[\bfseries\sffamily #1]}{\end{proof}}
\newenvironment{iclaim}[1][Claim]{\par\noindent\textbf{\sffamily #1:}\ }{\par}

%==================================================
% Shorthand Macros
%==================================================

\newcommand{\thm}[2]{\begin{Theorem}{#1}{}#2\end{Theorem}}
\newcommand{\cor}[2]{\begin{Corollary}{#1}{}#2\end{Corollary}}
\newcommand{\mlemma}[2]{\begin{Lemma}{#1}{}#2\end{Lemma}}
\newcommand{\mer}[2]{\begin{Exercise}{#1}{}#2\end{Exercise}}
\newcommand{\mprop}[2]{\begin{Prop}{#1}{}#2\end{Prop}}
\newcommand{\clm}[3]{\begin{claim}{#1}{#2}#3\end{claim}}
\newcommand{\wc}[2]{\begin{wconc}{#1}{}\setlength{\parindent}{1cm}#2\end{wconc}}
\newcommand{\thmcon}[1]{\begin{Theoremcon}#1\end{Theoremcon}}
\newcommand{\ex}[2]{\begin{Example}{#1}{}#2\end{Example}}
\newcommand{\dfn}[2]{\begin{Definition}{#1}{}#2\end{Definition}}
\newcommand{\dfnc}[2]{\begin{definition}{#1}{}#2\end{definition}}
\newcommand{\rmk}[2]{\begin{Remark}{#1}{}#2\end{Remark}}
\newcommand{\qs}[2]{\begin{question}{#1}{}#2\end{question}}
\newcommand{\pf}[2]{\begin{myproof}[#1]#2\end{myproof}}
\newcommand{\nt}[1]{\begin{note}#1\end{note}}
\newcommand{\iclm}[2]{\begin{iclaim}[#1]#2\end{iclaim}}
\newcommand{\mclm}[2]{\begin{myclaim}[#1]#2\end{myclaim}}

\newcommand{\cmtsolutionheading}{%
    \par\noindent\textbf{\textit{\sffamily Solution:}}\setlength{\parindent}{1cm}\par\smallskip\noindent%
}
\newenvironment{sol}{\cmtsolutionheading}{\par}
\newcommand{\solve}[1]{\cmtsolutionheading #1 \hfill\ensuremath{\blacksquare}}

\makeatother
